\centerline{\bf \Huge Геометрия > Алгебра}
\centerline{\bf \Huge }

\begin{flushright}
        \scriptsize{Что вершит судьбу алгебраиста в этом мире?\\ Некое незримое существо или геометр, \\подобно теореме косинусов парящей над системой?\\ По крайней мере истино то,\\ что алгебраист не властен даже над своим параметром...}
\end{flushright}

Наши друзья в этом листочке: \textbf{теорема косинусов, теорема синусов, уравнение окружности, прямой, отрезка, формула расстояния до прямой и расстояния между точками.}

\problem{1}
Докажите неравенство
$\sqrt{\left(x_1+x_2\right)^2+\left(y_1+y_2\right)^2} \leq \sqrt{x_1^2+y_1^2}+\sqrt{x_2^2+y_2^2}$.

В каком случае достигается равенство?

\problem{2}
Имеет ли решение в положительных числах система уравнений

$$\left\{
\begin{array}{l}
    x^2+x y+y^2=4 \\
    y^2+y z+z^2=9 \\
    z^2+z x+x^2=36 ?
\end{array}
\right.$$


\problem{3}
\textit{(Числа Грегори).} Докажите, что

$arctg\, 1+arctg\,\dfrac{1}{2} + arctg\, \dfrac{1}{3}=\dfrac{\pi}{2}.$

\problem{4}
\text { Решите систему }
$$\left\{\begin{array}{l}
    ay+b x=c, \\
    cx+a z=b, \\
    bz+c y=a .
\end{array}\right.$$

\problem{5}
Найдите все значения параметра $a$, при каждом из которых система

$$
\left\{\begin{array}{l}
y=|x-a|+a^2 \\
x^2+(y-2)^2=2
\end{array}\right.
$$


имеет единственное решение.

\newpage

\hspace{3mm}

\problem{6}
Эльзята выписала на доску все простые числа от 2 до 2027. За один ход ей можно заменить два числа $x, y$ на максимальное простое число, не превосходящее $\sqrt{x^2-xy+y^2}$. После какого-то количества ходов на доске останется одно число. Расул Владимирович обещал обещал Эльзяте заплатить столько тысяч рублей, сколько будет записано на доске в конце. Как Эльзяте обанкротить РВ?

\problem{7}
Найдите все значения параметра $a$, при каждом из которых система

$$
\left\{\begin{array}{l}
    2^x(y+1)\left(1-y \cdot 2^x\right)=a^3 \\
    \left(1+2^x\right)\left(1-y \cdot 2^x\right)=a
\end{array}\right.
$$

имеет хотя бы одно решение.
